\chapter{Conclusiones y trabajo futuro}
\label{cap:conclusiones}

El objetivo principal de este trabajo era diseñar y evaluar un sistema heterogéneo sobre la placa
Digilent Zybo, utilizando el procesador \ac{ARM} para controlar la aplicación y la
\ac{FPGA} para acelerar
una tarea de cálculo intensivo. Después del desarrollo realizado, se puede afirmar que este
objetivo se ha cumplido. El sistema final permite enviar una imagen desde el \ac{ARM},
transferirla a
la lógica programable mediante \ac{AXI} \ac{DMA}, ejecutar un modelo convolucional
generado con hls4ml y
recuperar los \textit{logits} de salida para obtener la clase predicha.

El camino hasta llegar a esa versión final no ha sido lineal. La principal dificultad no estuvo
solo en entrenar un modelo \ac{CIFAR}-10, sino en conseguir que dicho modelo cupiese
realmente en una
Zynq Z-7010. Muchas configuraciones que parecían razonables desde el punto de vista del número de
parámetros no eran viables después de pasar por Vivado, especialmente por el uso de
\ac{LUT} y por los
problemas de empaquetado de \textit{slices}. Esta parte del trabajo dejó claro que, en
una \ac{FPGA} pequeña, no
basta con hablar de tamaño del modelo: la forma en que las operaciones se sintetizan y se colocan
en el dispositivo es igual de importante.

\section{Cumplimiento de objetivos}

La primera parte del proyecto permitió poner en marcha la plataforma y entender el flujo completo
Vivado/Vitis. Se validó la ejecución \textit{bare-metal} sobre el \ac{ARM}, se configuró
la comunicación entre
el sistema de procesamiento y la lógica programable, y se trabajó con \ac{AXI}-Lite y AXI
\ac{DMA}. Estos
pasos fueron necesarios para que la parte de inteligencia artificial no quedase aislada
como un \ac{IP}
sin integración real.

La multiplicación de matrices de la fase 3 sirvió como punto intermedio. Fue una prueba más
controlada que la \ac{CNN}, pero suficiente para comprobar que Vitis \ac{HLS} podía generar hardware
competitivo y que el flujo \ac{ARM}--\ac{DMA}--\ac{FPGA} funcionaba correctamente.
Además, la comparación con una
implementación \ac{VHDL} ayudó a valorar el uso de \ac{HLS} con más criterio. La
herramienta no elimina la
necesidad de entender el hardware, pero sí reduce mucho el esfuerzo necesario para llegar a una
arquitectura paralela funcional.

En la fase 4 se consiguió generar e integrar un modelo \ac{CIFAR}-10 completo. La
precisión del modelo
se tuvo que equilibrar con los recursos disponibles en la Zybo. La versión final no es una red
grande, ni pretende competir con modelos modernos de visión artificial, pero sí cumple el propósito
del trabajo: demostrar una inferencia de \ac{IA} ejecutándose en hardware dedicado dentro
de una \ac{FPGA} de
recursos modestos.

Por último, la fase de evaluación permitió medir el beneficio de la aceleración. La inferencia en
\ac{FPGA} obtuvo un tiempo medio cercano a 2 ms, frente a unos 171 ms de la ejecución del
core \texttt{C++}
generado por hls4ml en el \ac{ARM}. Esto supone una mejora aproximada de 85 veces. Más
allá del número
exacto, lo relevante es que la aceleración se ha medido sobre la placa real e incluyendo el coste de
comunicación con el \ac{IP}.

\section{Lecciones aprendidas}

Una de las conclusiones más claras del proyecto es que \ac{HLS} facilita el desarrollo, pero no lo
convierte en automático. Las directivas, los tipos de datos y el factor de reutilización tienen un
impacto directo en los recursos finales. En varias ocasiones fue necesario volver al modelo,
modificar su arquitectura o ajustar la cuantización porque el diseño no entraba en la
\ac{FPGA}. En ese
sentido, el trabajo con hls4ml fue especialmente interesante: acerca mucho el mundo del aprendizaje
automático al diseño hardware, pero sigue exigiendo leer con cuidado los informes de síntesis,
utilización y temporización.

También ha sido importante distinguir entre los distintos niveles de validación. Al principio,
cuando la ejecución en placa devolvía resultados incoherentes, era tentador pensar en un fallo del
modelo. Sin embargo, parte del problema venía de cachés de Vivado/Vitis y de versiones antiguas del
\ac{IP} que seguían apareciendo en la integración.

Otro punto práctico ha sido la importancia de mantener los proyectos ordenados. En un flujo con
Vivado, Vitis, Vitis \ac{HLS}, modelos entrenados, \ac{IP} generados y \textit{scripts}
auxiliares, es fácil acabar
con carpetas repetidas o resultados antiguos mezclados con versiones nuevas. La separación entre
distintas fases y en cada una de ellas una carpeta con un proyecto Vivado/Vitis, hace que sea más
fácil de mantener y de explicar.

\section{Limitaciones}

La principal limitación del sistema es la propia placa. La Zybo Legacy es una plataforma muy útil
para aprender y prototipar, pero la Zynq Z-7010 tiene un margen reducido para implementar CNNs. El
modelo final tuvo que ser pequeño y muy ajustado en recursos. Esto afecta directamente a la
precisión que se puede esperar y limita la complejidad de las imágenes que puede clasificar con
robustez.

Otra limitación es que el \textit{benchmark} se ha centrado en latencia de inferencia sobre imágenes
concretas. Para una evaluación más completa sería necesario ejecutar un conjunto amplio de imágenes
de prueba y obtener métricas como precisión global, matriz de confusión o comportamiento por clase.
En este trabajo se priorizó cerrar el flujo completo hardware-software y medir la aceleración real
en placa.

También hay que tener cuidado al interpretar la confianza mostrada por terminal. El \ac{IP} devuelve
\textit{logits}, no probabilidades. La confianza se calcula en la aplicación con una
\textit{softmax} aproximada para
facilitar la lectura humana de los resultados. Por tanto, la clase predicha y los
\textit{logits} son más
importantes que el porcentaje exacto de confianza.

\section{Trabajo futuro}

Una línea clara de trabajo futuro sería mejorar el modelo manteniendo la restricción de recursos.
Para ello podrían explorarse técnicas como poda de pesos, cuantización consciente del entrenamiento
(\textit{quantization aware training}), arquitecturas separables en profundidad o una búsqueda más
sistemática de anchos de palabra por capa. El objetivo sería aumentar la precisión sin volver a
superar el límite de \ac{LUT} y \ac{BRAM} de la placa
\citep{han2015deep,howard2017mobilenetsefficientconvolutionalneural}.

También sería interesante estudiar mejor el uso de \ac{DSP}. En el diseño final el consumo de DSP es
bajo, mientras que las \ac{LUT} siguen siendo el recurso más comprometido. Un trabajo futuro podría
centrarse en forzar de forma más controlada el mapeo de multiplicaciones a \ac{DSP},
ajustar los anchos
de dato para favorecer ese mapeo y comparar distintas estrategias \ac{HLS}. No siempre será posible
trasladar lógica a \ac{DSP}, pero merece una exploración más profunda.

Otra mejora natural sería automatizar el flujo de generación. Actualmente el proceso implica
entrenar el modelo, exportarlo a hls4ml, generar el \ac{IP}, integrarlo en Vivado,
exportar el \ac{XSA} y
reconstruir la plataforma de Vitis. Parte de este flujo podría unificarse mediante
\textit{scripts} que
limpien versiones antiguas, regeneren el \ac{IP} y dejen preparada una carpeta final reproducible.
Esto reduciría errores de caché o de rutas, que han sido una fuente real de problemas durante el
proyecto.

Desde el punto de vista de la evaluación, el siguiente paso sería probar muchas más imágenes y no
solo casos individuales. Una aplicación que lea imágenes desde una tarjeta \ac{SD}, o
incluso desde una
cámara, permitiría evaluar el sistema de forma más parecida a un caso de uso real. También sería
interesante medir consumo energético, ya que una de las razones de usar \ac{FPGA} en
\textit{Edge-AI} no es solo
la latencia, sino la eficiencia energética \citep{sze2017efficientprocessingdeepneural}.

Por último, se podría comparar el acelerador \ac{FPGA} con una implementación \ac{ARM}
más optimizada, por
ejemplo usando NEON o librerías pensadas para microcontroladores y procesadores embebidos
\citep{lai2018cmsisnnefficientneuralnetwork}. Esto
permitiría situar mejor el resultado frente a una ruta software diseñada específicamente
para \ac{CPU}.
Aun así, el trabajo realizado ya muestra la idea principal: cuando el cálculo se adapta al hardware
y se integra correctamente con el sistema, la \ac{FPGA} puede reducir de forma notable la
latencia de
inferencia incluso en una placa de recursos limitados.
